% -*- coding: utf-8; -*-
\documentclass[english,submission]{programming}
%% First parameter: language is 'english'.
%% Second parameter: 'submission' for initial submission; remove for camera-ready (see §5.1).

\usepackage[backend=biber]{biblatex}
\addbibresource{bioscoop.bib}

%%
%% Additional packages (programming.cls already loads: amsmath, booktabs, listings,
%% xcolor, graphicx, hyperref, multirow, textcomp, and many others — do not reload them).
%%
\usepackage{csquotes}
\usepackage{amsthm}
\newtheorem{theorem}{Theorem}[section]

%% Custom colors for code listings
\definecolor{codegreen}{rgb}{0,0.6,0}
\definecolor{codegray}{rgb}{0.5,0.5,0.5}
\definecolor{codepurple}{rgb}{0.58,0,0.82}
\definecolor{backcolour}{rgb}{0.95,0.95,0.92}

%% Clojure listing language definition
\lstdefinelanguage{Clojure}{
  morekeywords={def, defn, defmacro, defrecord, defprotocol, defmethod,
                defmulti, let, if, cond, fn, ->, ->>, map, filter, reduce,
                when, case, for, binding, partial, apply, mapv, mapcat,
                instance?, assoc, update, update-in, into, vec, str},
  morecomment=[l]{;},
  morestring=[b]",
  sensitive=true
}

\lstdefinelanguage{BNF}{
  morecomment=[l]{;},
  morestring=[b]",
  sensitive=false
}

\lstset{
    language=Clojure,
    backgroundcolor=\color{backcolour},
    commentstyle=\color{codegreen},
    keywordstyle=\color{magenta},
    numberstyle=\tiny\color{codegray},
    stringstyle=\color{codepurple},
    basicstyle=\ttfamily\footnotesize,
    breakatwhitespace=false,
    breaklines=true,
    captionpos=t,
    keepspaces=true,
    numbers=left,
    numbersep=5pt,
    showspaces=false,
    showstringspaces=false,
    showtabs=false,
    tabsize=2,
    frame=single,
    escapeinside={(*@}{@*)}
}

%%
%% Paper classification details — REQUIRED for submission.
%% See §2.1 for perspective options: art, sciencetheoretical, scienceempirical, engineering.
%% See §2.2 for area options.
%%
\paperdetails{
  perspective=art,
  area={Interpreters, virtual machines and compilers},
  %% license=cc-by,  % default; change to cc-by-nc if needed
}

\begin{document}

\title{Bioscoop: A Domain-Specific Language for FFmpeg Filtergraph Composition}

\author{Daniel Szmulewicz}
\authorinfo{is an independent researcher working on programming languages,
  multimedia tooling, and creative coding. Contact him at
  \email{daniel.szmulewicz@gmail.com}.}
\affiliation{Independent Researcher}

\keywords{domain-specific languages, FFmpeg, filtergraph, multi-modal DSL, Clojure,
          bicameral parsing, intermediate representation, multi-stage programming}

%%
%% ACM CCS 2012 classification.
%% Generate the TeX code at https://dl.acm.org/ccs/ccs.cfm and paste it here.
%%
\begin{CCSXML}
<ccs2012>
<concept>
<concept_id>10011007.10011006.10011039.10011040.10011042</concept_id>
<concept_desc>Software and its engineering~Domain specific languages</concept_desc>
<concept_significance>500</concept_significance>
</concept>
<concept>
<concept_id>10011007.10011074.10011081.10011082</concept_id>
<concept_desc>Software and its engineering~Compilers</concept_desc>
<concept_significance>300</concept_significance>
</concept>
<concept>
<concept_id>10003752.10003809.10003636.10003646</concept_id>
<concept_desc>Theory of computation~Grammars and context-free languages</concept_desc>
<concept_significance>100</concept_significance>
</concept>
</ccs2012>
\end{CCSXML}

\ccsdesc[500]{Software and its engineering~Domain specific languages}
\ccsdesc[300]{Software and its engineering~Compilers}
\ccsdesc[100]{Theory of computation~Grammars and context-free languages}

\maketitle

%%
%% Abstract must follow the CIAKG+I structure:
%% Context, Inquiry, Approach, Knowledge, Grounding, Importance.
%%
\begin{abstract}
FFmpeg is the dominant open-source multimedia processing framework. Its filtergraph
syntax---a declarative language for composing processing pipelines---is expressive but
static: it offers no variables, no abstraction mechanisms, and no iteration. Filtergraphs of
any complexity must be written out in full or generated by string manipulation in the
host language.

We ask whether filtergraph generation can be treated as a \emph{compilation} problem
rather than an API problem, and what architectural properties such an approach entails.

We present Bioscoop, a domain-specific language compiler that targets FFmpeg filtergraph
syntax. Bioscoop provides two syntactically distinct front-ends---an external text-based DSL
parsed via GLL parsing, and an internal macro-based DSL embedded in Clojure---that
converge on an identical abstract syntax tree before evaluation.

We characterize several properties of the system. \emph{AST convergence}---the property
that both front-ends produce identical ASTs, enforced by a shared test suite---enables a
single definitional interpreter to handle both input paths, guaranteeing semantic equivalence
by construction. An \emph{isomorphic intermediate representation}, whose grammar is in
bijection with the FFmpeg filtergraph BNF grammar, reduces the back-end to a serializer
and enables round-trip correctness, established by structural induction. A \emph{first-class
iteration construct} enables data-driven generation of filtergraph topology and filter
arguments from computed values. \emph{Multi-stage programming} via \texttt{\&env} threads
Clojure lexical bindings into the DSL runtime. A \emph{fail-soft error discipline} accumulates
errors and returns typed sentinels rather than aborting, appropriate for interactive creative
coding. \emph{Dual resolution strategies}---runtime reflection on the JVM, static lookup in
GraalVM native image---maximize the affordances of each execution environment.

Bioscoop is implemented in Clojure and available as open-source software. Its formal
properties are established analytically and enforced by a test suite that runs identical
assertions against both compilation paths.

This work demonstrates that when a target language is already structured and compositional,
the intermediate representation can be isomorphic to the target grammar, the lowering
problem disappears, and the back-end reduces to a serializer. It situates Bioscoop within a
lineage of transpilers targeting domain notation languages---Sass targeting CSS, Hiccup
targeting HTML---where the value lies in providing the abstraction mechanisms that the
target language deliberately omits.
\end{abstract}

\section{Introduction}

FFmpeg is the dominant open-source multimedia processing framework, used for video and audio encoding, transcoding, filtering, and composition. Its filtergraph syntax---a declarative language for composing processing pipelines---is expressive but static: it offers no abstraction mechanisms, no variables, no iteration, and no conditional logic. Complex filtergraphs must be written out in full or generated by string manipulation outside the framework.

This paper describes Bioscoop, a compiler from a Lisp-like domain-specific language to FFmpeg filtergraph syntax. The key architectural decision that distinguishes Bioscoop from prior work is treating filtergraph generation as a \emph{compilation} problem rather than an API problem. Existing tools expose FFmpeg through function calls or builder patterns in the host language; the filtergraph remains implicit, assembled by the library as a string. Bioscoop instead defines a source language---with its own syntax, evaluation semantics, and abstraction mechanisms---that targets FFmpeg's filtergraph notation as its compilation output. To our knowledge, Bioscoop is the first source-to-source compiler whose target language is FFmpeg's filtergraph syntax.

Bioscoop addresses FFmpeg's expressiveness gap by providing:

\begin{itemize}
\item Named, reusable filter graph definitions (\texttt{defgraph})
\item Lexical binding (\texttt{let})
\item A first-class iteration construct (\texttt{for}) that generates parallel filtergraph structures
\item Parameterized filter arguments computed at compile time
\item Spec-driven parameter validation
\end{itemize}

The system supports two syntactically distinct front-ends that share a single evaluation pipeline. We describe the architectural consequences of this design, with particular attention to the properties that distinguish Bioscoop from string-generation approaches and from prior DSL work in the multimedia domain.

\subsection{Motivating Example}

Consider generating a slideshow that cross-fades $n$ images. In raw FFmpeg, every intermediate label must be written explicitly. For ten images this requires approximately ninety characters of filtergraph syntax, all of which must be manually regenerated if $n$ changes. In Bioscoop:

\begin{lstlisting}[caption={Parameterized slideshow in Bioscoop}]
(defn slideshow [n]
  (bioscoop
    (for [i (range n)]
      [[(str i)] (loop {:loop 124 :size 1})
                 [(str "l" i)]])
    (for [i (range 1 n)]
      [[(if (= i 1) "l0" (str "v" i))
        (str "l" i)]
       (xfade {:transition "fade"
               :duration   1
               :offset     (+ (* i 4) 3)})
       [(str "v" i)]])))
\end{lstlisting}

\texttt{(slideshow 10)} produces the equivalent of ninety lines of hand-written FFmpeg filtergraph syntax. Changing $n$, the transition type, or the timing requires changing a single argument.

\section{Background and Related Work}

\subsection{FFmpeg Filtergraph Syntax}

An FFmpeg filtergraph is a semicolon-separated sequence of filterchains, each a comma-separated sequence of filters. Filters carry optional bracketed input and output link labels. The BNF grammar is:

\begin{lstlisting}[language=BNF, caption={FFmpeg filtergraph grammar (simplified)}]
filtergraph = filterchain (";" filterchain)*
filterchain = filter ("," filter)*
filter      = linklabel* filter-spec linklabel*
filter-spec = name ("=" arguments)?
linklabel   = "[" name "]"
\end{lstlisting}

The grammar is context-free and unambiguous. There are no variables, no abstraction mechanisms, and no iteration. Filtergraphs of any complexity must be written as single expressions.

\subsection{Existing Approaches}

Existing tools for programmatic FFmpeg filtergraph generation fall into three categories. \emph{API wrappers} such as ffmpeg-python~\cite{ffmpegpython} and fluent-ffmpeg~\cite{fluentffmpeg} expose FFmpeg's options as function calls but treat the filtergraph itself as an opaque string to be concatenated. \emph{Higher-level abstractions} such as MoviePy~\cite{moviepy} hide FFmpeg almost entirely, trading expressiveness for simplicity.

A third category is represented by python-ffmpegio~\cite{ffmpegio}, which provides a structured Python object model (\texttt{Filter}, \texttt{Chain}, \texttt{Graph}) that serializes to filtergraph strings and can parse native filtergraph strings back into the same representation. This is the closest prior work to Bioscoop's intermediate representation. However, python-ffmpegio's filtergraph builder is an API, not a source language: there is no separate syntax, no named definitions with lexical scope, no iteration construct, and no parameter validation. Users construct filtergraphs by combining Python objects with operator overloading (\texttt{+}, \texttt{|}). Bioscoop differs in kind: it provides a source language with its own parser and evaluator, named abstraction mechanisms that are compiled away before serialization, and spec-driven validation at compile time.

\subsection{DSL Implementation in Lisp}

The design of internal DSLs in Lisp is well-studied. Hudak's taxonomy~\cite{hudak} distinguishes shallow embeddings, where host language constructs carry DSL semantics, from deep embeddings, where DSL programs are represented as data and interpreted by a separate evaluator. Bioscoop uses a deep embedding: Clojure forms are rewritten into a tagged vector AST that is then evaluated by a definitional interpreter.

Racket's \texttt{\#lang} mechanism~\cite{racket} and Clojure's macro system both support what we call AST convergence naturally, though we are not aware of prior work that names or formalizes this property in the context of multi-modal DSLs.

\section{Architecture}

\subsection{Overview}

Bioscoop consists of six layers:

\begin{enumerate}
\item \emph{Two front-ends}: a GLL parser for the external text DSL and a macro for the internal Clojure DSL
\item \emph{A shared AST}: tagged Clojure vectors produced identically by both front-ends
\item \emph{A definitional interpreter}: \texttt{transform-ast}, a multimethod dispatching on AST node type
\item \emph{An intermediate representation}: algebraic records (\texttt{Filter}, \texttt{FilterChain}, \texttt{FilterGraph})
\item \emph{A back-end}: \texttt{to-ffmpeg}, a serializer from IR to filtergraph string
\item \emph{An error channel}: a dynamically-scoped accumulator for compile-time errors
\end{enumerate}

\subsection{AST Convergence}

Both front-ends produce ASTs consisting of tagged Clojure vectors. The node types are: \texttt{:program}, \texttt{:compose}, \texttt{:for-binding}, \texttt{:let-binding}, \texttt{:binding}, \texttt{:padded-graph}, \texttt{:list}, \texttt{:map}, \texttt{:symbol}, \texttt{:keyword}, \texttt{:string}, \texttt{:number}, \texttt{:boolean}, \texttt{:label}, \texttt{:graph-definition}.

Formally, let $P_e : L_e \rightarrow AST$ be the external parser and $M_i : L_i \rightarrow AST$ be the macro. AST convergence holds when:

\begin{equation}
\forall p \in L_e \cap L_i : P_e(p) = M_i(p)
\end{equation}

where $L_e \cap L_i$ denotes programs expressible in both surface syntaxes. This enables a single transformer $T : AST \rightarrow IR$ to handle both paths, guaranteeing semantic equivalence by construction.

What makes this achievable is what Krishnamurthi~\cite{krishnamurthi} calls a \emph{bicameral} parsing pipeline. In Clojure, the reader stage converts source text into standard Clojure data structures---lists, vectors, symbols, maps---and the macro receives these directly at compile time. Crucially, Clojure's reader output \emph{is} standard Clojure data: not an intermediate format, not special syntax objects as in Racket, but the same lists and vectors manipulated anywhere in a Clojure program. The \texttt{form->ast} function can therefore traverse a Clojure form using ordinary Clojure functions (\texttt{cond}, \texttt{first}, \texttt{mapv}) and produce tagged vectors.

Bicamerality makes AST convergence \emph{possible}; it does not make it \emph{guaranteed}. Any two independent implementations of a source-to-AST transformation will diverge unless actively maintained. The term \emph{homoiconicity}, commonly invoked to explain this kind of macro power, is imprecise: as Krishnamurthi demonstrates, it describes either a trivially universal property (any language with strings can represent programs as data) or a specific and problematic feature (\texttt{eval}), neither of which is the actual mechanism here. The mechanism is the bicameral pipeline. What makes the convergence guarantee hold is a deliberate design decision: \texttt{form->ast} was written to emit the same tagged-vector node types that the GLL parser emits, and the test suite enforces this pairing at every new construct.

The external front-end uses Instaparse~\cite{instaparse}, which implements the GLL (Generalized LL) parsing algorithm~\cite{gll}. GLL handles the full class of context-free grammars, including ambiguous and left-recursive ones, while retaining the structural clarity of recursive descent. For an unambiguous grammar, as Bioscoop's is, GLL runs in linear time.

The internal front-end is the \texttt{bioscoop} macro:

\begin{lstlisting}[caption={The bioscoop macro}]
(defmacro bioscoop [& forms]
  (let [ast-nodes   (mapv form->ast forms)
        program-ast (vec (concat [:program] ast-nodes))
        locals      (keys &env)]
    `(binding [config/*dynamic-resolution* true]
       (let [env# (reduce (fn [e# [k# v#]]
                            (env-put e# k# v#))
                          (make-env)
                          ~(mapv (fn [sym]
                                   [(str sym) sym])
                                 locals))]
         (run-ast ~program-ast env#)))))
\end{lstlisting}

\texttt{form->ast} is a structural transformation from Clojure forms to AST nodes. It handles all surface syntax cases: function calls, symbols, keywords, strings, numbers, booleans, padded graphs, maps, and the \texttt{for}, \texttt{let}, and \texttt{compose} special forms.

\subsection{Multi-Stage Programming via \texttt{\&env}}

A significant design challenge is threading Clojure lexical bindings into the DSL's runtime environment. Consider:

\begin{lstlisting}[caption={Parameterized iteration with lexical scope}]
(defn n-transition [n offset]
  (bioscoop
    (for [i (range n)]
      [[(str "v" i)]
       (xfade {:transition "fade"
               :duration   1
               :offset     (+ i offset (* i offset))})
       [(str "t" (inc i))]])))
\end{lstlisting}

Here \texttt{n} and \texttt{offset} are Clojure locals, not DSL bindings. The \texttt{bioscoop} macro must make them visible to the DSL evaluator at runtime.

The solution uses multi-stage programming~\cite{taha}: \texttt{(keys \&env)} captures the names of all local bindings at macro \emph{expansion time}, as a sequence of symbols. The macro then generates code that, at \emph{runtime}, evaluates those symbols to their current values and injects them into the DSL environment:

\begin{lstlisting}[caption={The locals injection mechanism}]
; At expansion time, for locals [n offset]:
; (keys &env) = (n offset)
; (mapv (fn [sym] [(str sym) sym]) locals)
; produces the literal: [["n" n] ["offset" offset]]
;
; At runtime, n=10 and offset=6, so:
; [["n" 10] ["offset" 6]]
; is injected into the DSL env
\end{lstlisting}

The key asymmetry is that \texttt{(str sym)} converts the symbol to its string name at expansion time---the only moment the name is available as a symbol---while \texttt{sym} remains unevaluated so that the generated code reads its value from the live call frame at runtime. This is a two-phase computation: names extracted at compile time, values read at runtime.

\section{The Intermediate Representation}

\subsection{Algebraic Structure}

The IR is defined by three record types:

\begin{lstlisting}[caption={IR record definitions}]
(defrecord Filter      [name args
                        input-labels output-labels])
(defrecord FilterChain [filters])
(defrecord FilterGraph [chains])
\end{lstlisting}

These define a context-free grammar over Clojure records:

\begin{align*}
\text{FilterGraph} &::= [\text{FilterChain}^*] \\
\text{FilterChain}  &::= [\text{Filter}^+] \\
\text{Filter}       &::= \text{name} \times \text{args} \times \text{labels}^* \times \text{labels}^*
\end{align*}

Well-formedness is enforced by construction: \texttt{make-filtergraph} accepts only vectors of \texttt{FilterChain} records, \texttt{make-filterchain} accepts only vectors of \texttt{Filter} records. Invalid structures cannot be constructed through the public API.

\subsection{Isomorphism with the Target Grammar}

The IR grammar and the FFmpeg filtergraph BNF grammar are in bijection:

\begin{table}
\centering
\caption{IR to FFmpeg Grammar Correspondence}
\label{tab:isomorphism}
\begin{tabular}{ll}
\toprule
\textbf{IR} & \textbf{FFmpeg BNF} \\
\midrule
\texttt{FilterGraph} & \texttt{filtergraph} \\
\texttt{FilterChain} & \texttt{filterchain} \\
\texttt{Filter} & \texttt{filter} \\
\texttt{:name} & \texttt{filter-name} \\
\texttt{:args} & \texttt{filter-arguments} \\
\texttt{:input-labels} & \texttt{input-linklabels} \\
\texttt{:output-labels} & \texttt{output-linklabels} \\
\texttt{;} between chains & \texttt{filterchain} separator \\
\texttt{,} between filters & \texttt{filter} separator \\
\bottomrule
\end{tabular}
\end{table}

This isomorphism has a significant architectural consequence: the back-end \texttt{to-ffmpeg} is a \emph{serializer} rather than a compiler pass. It traverses a structure already in the shape of the target and emits delimiters:

\begin{lstlisting}[caption={to-ffmpeg as serializer}]
(extend-protocol Renderable
  Filter
  (to-ffmpeg [f]
    (let [in  (str/join ""
                (map #(str "[" % "]")
                     (:input-labels f)))
          out (str/join ""
                (map #(str "[" % "]")
                     (:output-labels f)))
          args (when (:args f)
                 (str "=" (str/join ":"
                   (map (fn [[k v]]
                          (str (name k) "=" v))
                        (:args f)))))]
      (str in (:name f) args out)))
  FilterChain
  (to-ffmpeg [fc]
    (str/join "," (map to-ffmpeg (:filters fc))))
  FilterGraph
  (to-ffmpeg [fg]
    (str/join ";" (map to-ffmpeg (:chains fg)))))
\end{lstlisting}

No information is added or removed. The serializer is a homomorphism from the IR algebra to strings under concatenation.

\subsection{Round-Trip Correctness}

The isomorphism also enables a bidirectional parser: \texttt{ffmpeg-parser} parses native FFmpeg filtergraph strings into the same record structure that the DSL compiler produces. We state the round-trip property formally and give a proof sketch by structural induction.

Let $\mathcal{S}$ denote the set of well-formed FFmpeg filtergraph strings, $\mathcal{I}$ the set of well-formed IR values, $\sigma : \mathcal{I} \rightarrow \mathcal{S}$ the serializer \texttt{to-ffmpeg}, and $\pi : \mathcal{S} \rightarrow \mathcal{I}$ the parser \texttt{ffmpeg-parser}.

\begin{theorem}[Round-trip correctness]
For all $v \in \mathcal{I}$, $\pi(\sigma(v)) = v$.
\end{theorem}

\begin{proof}[Proof sketch]
By structural induction on the IR grammar.

\textbf{Base case: \texttt{Filter}.} Let $f = \langle n, a, \ell_i, \ell_o \rangle$ where $n$ is the filter name, $a$ the argument map, $\ell_i$ the input labels, and $\ell_o$ the output labels. The serializer produces:
\[
\sigma(f) = \underbrace{[\ell_{i_1}]\cdots[\ell_{i_k}]}_{\text{input labels}} \; n \; \underbrace{=a_1:a_2:\cdots}_{\text{args}} \; \underbrace{[\ell_{o_1}]\cdots[\ell_{o_m}]}_{\text{output labels}}
\]
The FFmpeg grammar rule \texttt{filter = linklabel* filter-spec linklabel*} parses this string unambiguously, recovering $n$ from \texttt{filter-name}, $a$ from \texttt{filter-arguments}, $\ell_i$ from leading \texttt{linklabels}, and $\ell_o$ from trailing \texttt{linklabels}. Since labels are stored as first-class fields (not metadata), equality holds: $\pi(\sigma(f)) = f$.

\textbf{Inductive case: \texttt{FilterChain}.} Let $c = [f_1, \ldots, f_n]$ where each $f_j$ is a \texttt{Filter}. The serializer produces $\sigma(f_1) \mathbin{,} \sigma(f_2) \mathbin{,} \cdots \mathbin{,} \sigma(f_n)$. The FFmpeg grammar rule \texttt{filterchain = filter ("," filter)*} partitions this string at each top-level comma (commas do not appear in filter names, argument keys, or labels), yielding $n$ substrings. By the base case, $\pi(\sigma(f_j)) = f_j$ for each $j$. Therefore $\pi(\sigma(c)) = [f_1, \ldots, f_n] = c$.

\textbf{Inductive case: \texttt{FilterGraph}.} Let $g = [c_1, \ldots, c_m]$ where each $c_k$ is a \texttt{FilterChain}. The serializer produces $\sigma(c_1) \mathbin{;} \sigma(c_2) \mathbin{;} \cdots \mathbin{;} \sigma(c_m)$. The FFmpeg grammar rule \texttt{filtergraph = filterchain (";" filterchain)*} partitions this string at each top-level semicolon (semicolons do not appear within filterchains). By the inductive case, $\pi(\sigma(c_k)) = c_k$ for each $k$. Therefore $\pi(\sigma(g)) = [c_1, \ldots, c_m] = g$.
\end{proof}

The converse---$\sigma(\pi(s)) = s$ for all $s \in \mathcal{S}$---does not hold in general because the FFmpeg serialization format admits multiple syntactic representations of the same filtergraph (equivalent argument orderings, optional whitespace). Bioscoop normalizes these to a canonical form during parsing. The round-trip that is guaranteed is $\text{IR} \rightarrow \text{string} \rightarrow \text{IR}$, which is the direction that matters for the compiler's correctness.

\subsection{Labels as First-Class Fields}

An earlier version of Bioscoop stored labels as Clojure metadata on \texttt{Filter} records. This was architecturally incorrect: Clojure's specification states that metadata does not affect equality, so two filters with different labels were equal. This violated the IR grammar, where \texttt{[in]scale} and \texttt{scale} are distinct constructs.

The current design stores labels as explicit fields. Two filters with different labels are correctly unequal. The \texttt{Sinkable} protocol, which existed solely to manage the metadata workaround, is eliminated. The IR now enforces the grammar's semantics through the type system.

\section{The Evaluator}

\subsection{Definitional Interpreter in Environment-Passing Style}

\texttt{transform-ast} is a definitional interpreter~\cite{reynolds} written in environment-passing style. Each \texttt{defmethod} gives the meaning of one syntactic form by structural recursion over the parse tree:

\begin{lstlisting}[caption={The evaluator entry point}]
(defmulti transform-ast*
  (fn [node env] (first node)))

(defn transform-ast [node env]
  (trace> node env (transform-ast* node env)))
\end{lstlisting}

The evaluator is:

\begin{itemize}
\item \emph{Single-pass}: one recursive descent over the AST
\item \emph{Eagerly evaluated}: applicative-order evaluation
\item \emph{Lexically scoped}: the environment is a chain of frames linked by parent pointers
\item \emph{Total}: every well-formed input produces a \texttt{FilterGraph}, possibly empty
\end{itemize}

Tracing is delegated to \texttt{trace>}, defined in a separate \texttt{bioscoop.trace} namespace:

\begin{lstlisting}[caption={The trace> helper}]
(defn trace> [node env result]
  (tap> (cond-> {:node-type (first node)
                 :node      node
                 :env       env
                 :result    result}
          (instance? FilterGraph result)
          (assoc :ffmpeg (to-ffmpeg result))))
  result)
\end{lstlisting}

Every node's type, environment state, and result are forwarded to Clojure's \texttt{tap>} facility. The \texttt{:ffmpeg} key is added conditionally via \texttt{cond->}: it is only present when \texttt{result} is a \texttt{FilterGraph}, so the tap map never carries a \texttt{nil} serialization. Absence of \texttt{:ffmpeg} is itself the signal that the result is a partial value---no separate boolean flag is needed. Callers register tap subscribers to capture trace data; no subscribers means no overhead.

\subsection{Environment Model}

The environment is a Clojure map extended with a parent pointer:

\begin{lstlisting}[caption={Environment operations}]
(defn make-env
  ([]       {:errors (atom [])})
  ([parent] (assoc {:errors (:errors parent)} :parent parent)))

(defn env-get [env sym]
  (if (contains? env sym)
    (get env sym)
    (when-let [parent (:parent env)]
      (env-get parent sym))))

(defn env-put [env sym val]
  (assoc env sym val))
\end{lstlisting}

New scopes are created with \texttt{(make-env parent)}, establishing the parent chain for lexical scope traversal; child scopes share the parent's error atom, so errors raised at any depth surface in the single atom inspected after compilation.

\texttt{contains?} rather than \texttt{if-let} is used to correctly handle bindings to falsy values. New scopes are created with \texttt{(make-env parent)}, establishing the parent chain for lexical scope traversal.

\subsection{The \texttt{for} Construct}

The \texttt{for} construct is the most significant addition to the filtergraph language. It iterates over any seqable Clojure value and produces a \texttt{FilterGraph} by merging the results of evaluating the body expression for each iteration value:

\begin{lstlisting}[caption={for-binding evaluation}]
(defmethod transform-ast* :for-binding
  [[_ [_ sym-name] range-node body] env]
  (let [xs (transform-ast range-node env)]
    (vec (for [x xs
               :let [loop-env (env-put env sym-name x)]]
           (transform-ast body loop-env)))))
\end{lstlisting}

The method returns a plain Clojure vector of per-iteration results. Composition into a \texttt{FilterGraph} is handled upstream by \texttt{promote-to-filtergraph}, which contains a \texttt{sequential?} branch that recursively promotes each element and combines them with \texttt{compose-filtergraphs}:

\begin{lstlisting}[caption={promote-to-filtergraph* with sequential? branch}]
(defn promote-to-filtergraph* [f]
  (fn promote [x env]
    (cond
      (instance? FilterGraph x) x
      (instance? FilterChain x) (make-filtergraph [x])
      (instance? Filter x)
        (make-filtergraph [(make-filterchain [x])])
      (sequential? x)
        (apply compose-filtergraphs
               (map #(promote % env) x))
      :else (f env x :not-a-filtergraph))))
\end{lstlisting}

This separation of concerns means the \texttt{for-binding} evaluator need not know whether its callers want a \texttt{FilterGraph}, a label list, or some other interpretation. The range expression is evaluated against the current environment, so it can reference DSL bindings and Clojure locals injected via \texttt{\&env}. All of Clojure's sequence-producing functions---\texttt{range}, \texttt{map}, \texttt{filter}, \texttt{iterate}---are available as range expressions.

The \texttt{for} construct also appears in label position, where it generates a sequence of label strings rather than a filtergraph:

\begin{lstlisting}[caption={for in label position}]
; Generates input labels ["b0c" "b1c" ... "b(n-1)c"]
; for the stacking filter:
[[(for [i (range n)] (str "b" i "c"))] stacking]
\end{lstlisting}

This is handled naturally by the padded-graph evaluator: when a label expression evaluates to a sequential value, the handler flattens it via \texttt{mapcat}. Since \texttt{for-binding} returns a plain vector, a \texttt{for} expression in label position produces the expected sequence of label strings without any special dispatch on the node type.

\subsection{Chain Flattening}

The \texttt{chain} construct flattens heterogeneous sequences of \texttt{Filter}, \texttt{FilterChain}, and single-chain \texttt{FilterGraph} values into a single \texttt{FilterChain}:

\begin{lstlisting}[caption={chain flattening via promote-to-filterchain}]
"chain"
(make-filterchain
  (vec (mapcat #(promote-to-filterchain % env)
               transformed-args)))
\end{lstlisting}

\texttt{promote-to-filterchain} is a normalization function parameterized by an error handler:

\begin{lstlisting}[caption={promote-to-filterchain* factory}]
(defn promote-to-filterchain* [f]
  (fn [x env]
    (cond
      (instance? Filter x)      [x]
      (instance? FilterChain x) (:filters x)
      (instance? FilterGraph x)
        (if (= 1 (count (:chains x)))
          (:filters (first (:chains x)))
          (f env x :chain-parallel-filtergraph))
      :else (f env x :not-a-filtergraph))))
\end{lstlisting}

This allows \texttt{defgraph}-defined named graphs to be used transparently as chain stages, since a \texttt{defgraph} with a linear body produces a single-chain \texttt{FilterGraph} that flattens into its constituent filters. Multi-chain filtergraphs---representing parallel structure---are rejected as a type error.

\section{Error Handling}

\subsection{Fail-Soft with Typed Sentinel}

Bioscoop targets interactive creative coding at the REPL. In this context, a hard-failure model---where compilation errors throw exceptions and abort evaluation---is counterproductive. Partial results and immediate feedback matter more than strict failure boundaries.

The error discipline is fail-soft with explicit sentinel: all error paths return \texttt{(make-filtergraph [])} and accumulate a structured error record in a dynamically-scoped error atom. The first error is printed immediately at the point of accumulation. The atom is inspectable after compilation completes:

\begin{lstlisting}[caption={Error accumulation with deduplication}]
(defn accumulate-error*
  ([env error]
   (accumulate-error* (make-filtergraph []) env error))
  ([return-val env error]
   (let [info (ex-data error)]
     (when-not (some #(= (ex-data %) info) @(:errors env))
       (log/warn (if *warn-verbose* info (:error-type info)))
       (swap! (:errors env) conj error)))
   return-val))
\end{lstlisting}

The \texttt{when-not} guard prevents the same error from being logged and accumulated more than once. This matters when error-producing nodes appear in paths that are traversed multiple times---for instance, a graph definition that is referenced repeatedly. Without the guard, the error atom would accumulate duplicate entries and the log would be flooded.

The \texttt{return-val} parameter, defaulting to an empty \texttt{FilterGraph}, allows error handlers to return type-appropriate sentinels. Filterchain contexts return \texttt{[]} since their results feed \texttt{mapcat}. The promotion functions are wired directly to the accumulator:

\begin{lstlisting}[caption={Direct wiring of promotion functions to error handler}]
(def promote-to-filtergraph
  (promote-to-filtergraph* accumulate-error))

(def promote-to-filterchain
  (promote-to-filterchain* (partial accumulate-error [])))
\end{lstlisting}

\texttt{promote-to-filtergraph} passes \texttt{accumulate-error} directly, since the factory will supply the empty-\texttt{FilterGraph} sentinel via the default arity. \texttt{promote-to-filterchain} pre-applies the \texttt{[]} sentinel inline, since filterchain contexts must return an empty vector for \texttt{mapcat} rather than an empty \texttt{FilterGraph}.

\subsection{Invariant: transform-ast Always Returns FilterGraph}

The central invariant of the evaluator is that \texttt{transform-ast}, when called in filter-producing position, always returns a \texttt{FilterGraph}. Error paths return an empty \texttt{FilterGraph} rather than a map, an exception, or \texttt{nil}. This eliminates defensive type-checking throughout the evaluator: no method needs to guard against receiving a non-\texttt{FilterGraph} from a recursive call.

The invariant is established at the program boundary by \texttt{promote-to-filtergraph}, which is the only function that transitions from the untyped world of evaluated expressions to the typed world of IR values. Every method that consumes a recursive result calls \texttt{promote-to-filtergraph} before working with it.

\section{Spec-Driven Validation}

Filter parameters are validated against \texttt{clojure.spec}~\cite{clojurespec} specifications before construction. Each filter has a corresponding spec defining the valid types and ranges of its parameters:

\begin{lstlisting}[caption={Example spec for the lagfun filter}]
(s/def ::decay  (s/double-in :min 0 :max 1))
(s/def ::planes (s/int-in 0 16))
(s/def ::lagfun (s/keys :opt-un [::decay ::planes]))
\end{lstlisting}

Parameter maps use unqualified keys in DSL source. \texttt{spec-aware-namespace-map} resolves these to their correct namespace-qualified forms by consulting the spec registry:

\begin{lstlisting}[caption={Namespace resolution via spec introspection}]
(defn spec-aware-namespace-map [spec-keyword m]
  (let [spec-form  (s/form spec-keyword)
        key-mapping
          (into {}
            (map (fn [qk]
                   [(keyword (name qk)) qk])
                 (all-keys spec-form)))]
    (into {}
      (map (fn [[k v]]
             [(get key-mapping k k)
              (coerce-numeric v)])
           m))))
\end{lstlisting}

The \texttt{coerce-numeric} function converts Clojure \texttt{Ratio} values (produced by integer division) to \texttt{Double}, ensuring compatibility with specs defined via \texttt{s/double-in}. This is a latent type error that expression arguments expose: \texttt{(/ 1 3)} produces a \texttt{Ratio} in Clojure, not a \texttt{Double}.

\section{Discussion}

\subsection{Applicability of the Isomorphic IR Approach}

The absence of a lowering problem in Bioscoop is a consequence of domain analysis, not a limitation of ambition. FFmpeg filtergraphs are already a high-level, compositional, declarative language. The source language and target language are at the same level of abstraction. There is no semantic gap to bridge, only a syntactic and ergonomic one.

This situation is not universal. A compiler from a higher-level language to machine code must lower closures, garbage collection, type erasure, and many other abstractions that have no counterpart in the target. The isomorphic IR approach is appropriate when:

\begin{enumerate}
\item The target language is already structured and compositional
\item The source language adds abstraction mechanisms rather than new semantic concepts
\item The primary function of the compiler is elaboration---name resolution, validation, iteration expansion---rather than semantic transformation
\end{enumerate}

FFmpeg filtergraphs satisfy all three conditions. Configuration languages, build system descriptors, query languages, and other declarative domain notations typically do as well.

\subsection{AST Convergence as a Design Pattern}

AST convergence is a design decision, not a consequence of Clojure's bicameral parsing pipeline alone. Without deliberate maintenance, the two front-ends will diverge as the language evolves: new constructs will be added to the macro path but not the grammar, or vice versa, and the behavioral equivalence guarantee will silently erode.

The mechanism that enforces convergence is the shared AST node type set. Both front-ends must produce nodes from the same enumerated set. When a new construct is added---\texttt{for}, for instance---it must be added to both the grammar and \texttt{form->ast} simultaneously. The test suite enforces this by running identical assertions against both compilation paths:

\begin{lstlisting}[caption={Convergence test pattern}]
(deftest macro-string-convergence
  (is (= (to-ffmpeg
           (compile-dsl
             "(chain (scale 1920 1080) (eq))"))
         (to-ffmpeg
           (bioscoop
             (chain (scale 1920 1080) (eq)))))))
\end{lstlisting}

\subsection{Single-Pass Architecture}

The evaluator is single-pass in the sense that one recursive descent over the AST produces the final IR. There are no analysis phases, no optimization passes, and no separate code generation step. The \texttt{:program} method processes graph definitions before regular expressions, threading the environment through each definition in turn so that a later definition may reference an earlier one; this establishes an ordering constraint, but it does not constitute a second pass: no new representation is produced, and every graph definition is resolved through the same environment-passing mechanism as the rest of the evaluator.

This characterization aligns with the classical definition: a compiler pass is a complete traversal of the program representation that produces a new representation. Since Bioscoop uses one representation throughout, it has one pass regardless of the traversal order of individual nodes.

\subsection{Dual Resolution Strategies}

Bioscoop targets two distinct execution environments: the JVM Clojure runtime and GraalVM native image. The \texttt{resolve-function} multimethod selects between two resolution strategies via the \texttt{*dynamic-resolution*} dynamic var.

The JVM path uses runtime reflection via \texttt{ns-resolve}. This affords capabilities unique to the JVM's open-world execution model: vars defined after the program started, graphs interned via \texttt{defgraph} in arbitrary namespaces, and fully qualified symbol resolution across dynamically loaded code. User-defined filtergraph functions are resolved by searching the current namespace and its transitive dependencies at runtime, enabling interactive development patterns where graphs are defined and redefined at the REPL without restarting the system.

The GraalVM native image path uses a static lookup table populated at compile time from the public vars of \texttt{bioscoop.built-in} and \texttt{clojure.core}. GraalVM's closed-world compilation model restricts dynamic class loading and requires explicit reflection configuration. The static table is the correct response to these constraints: it is fast, allocation-free, and requires no reflection at runtime.

These are not two implementations of the same mechanism but two distinct strategies, each maximizing the affordances of its execution environment---for graph resolution as much as for filter-function resolution. The \texttt{*dynamic-resolution*} dynamic var is the abstraction boundary between them, selected automatically by the macro path and configurable for programmatic use.

\subsection{Limitations}

The current system has one notable limitation.

\emph{Label arity at compose boundaries}: The \texttt{for} construct can generate $n$ parallel filtergraph chains, each with computed output labels. But a subsequent \texttt{compose} stage that collects all $n$ outputs as inputs to a single filter---such as \texttt{concat} or \texttt{xstack} with $n$ inputs---must still list those labels explicitly. A splice operator for label sequences at compose boundaries would close this gap.

\section{Related Work}

\subsection{Language Workbenches}

JetBrains MPS~\cite{mps} and Xtext~\cite{xtext} provide tooling for multi-modal DSL development, typically maintaining separate grammars for different surface syntaxes. Neither targets the problem of shared intermediate representation between modalities.

\subsection{Nanopass Framework}

The Nanopass Framework~\cite{nanopass} for Scheme decomposes compilation into many small, explicitly named passes, each operating on a distinct intermediate language. This approach maximizes comprehensibility of individual passes at the cost of managing many intermediate representations. Bioscoop's single-representation approach is at the opposite end of this design spectrum, appropriate because the source and target languages are at the same level of abstraction.

\subsection{Typed DSLs}

Typed embeddings of DSLs in functional languages~\cite{carette} use the host type system to enforce DSL well-formedness statically. Bioscoop uses runtime validation via \texttt{clojure.spec} rather than static types, trading compile-time guarantees for flexibility and REPL interactivity.

\subsection{Source-to-Source Compilation for Domain Notation Languages}

The pattern of targeting an existing domain notation language as a compilation output appears in several successful systems. Sass and Less~\cite{sass} target CSS, adding variables, nesting, mixins, and functions to a language that is deliberately simple and declarative. These systems demonstrate that the source language / notation language split is an architecturally stable arrangement: CSS has gained custom properties and other features, yet Sass remains in wide use because a compilation model can offer abstraction capabilities that a notation language should deliberately omit.

In the Clojure ecosystem, Hiccup~\cite{hiccup} provides a DSL---represented as Clojure vectors---that targets HTML. The structural parallel with Bioscoop is real: both generate a well-formed textual notation from a tree-structured source. The analogy is partial, however, because HTML is markup rather than a processing language and its grammar is permissive rather than compositionally constrained. FFmpeg's filtergraph syntax is closer to CSS in this respect: it is a structured, compositional notation for a processing pipeline, with a well-defined grammar and precise semantics for filter composition.

Bioscoop fits this pattern. Its target language is a public, stable API---\texttt{libavfilter}'s filtergraph syntax---that is intentionally devoid of abstraction mechanisms. The compilation model is stable for the same reason Sass is stable targeting CSS: the target language's simplicity is a deliberate design decision, not a temporary deficiency to be corrected upstream.

\section{Conclusion}

We have described Bioscoop, a domain-specific language compiler for FFmpeg filtergraph composition. Its primary contributions are:

\begin{itemize}
\item \emph{Novel compilation target}: to our knowledge, the first source-to-source compiler whose target language is FFmpeg's filtergraph syntax, as distinct from API wrappers and builder libraries that construct filtergraphs programmatically without a source language
\item \emph{AST convergence}: a deliberately maintained architectural property ensuring that two front-ends produce identical ASTs before evaluation, guaranteeing semantic equivalence by construction
\item \emph{Isomorphic IR}: an intermediate representation whose grammar is in bijection with the target language's grammar, reducing the back-end to a serializer
\item \emph{First-class iteration}: a \texttt{for} construct that enables data-driven generation of filtergraph structure---the first such abstraction mechanism in FFmpeg tooling
\item \emph{Multi-stage locals injection}: a mechanism using \texttt{\&env} to thread Clojure lexical bindings into the DSL runtime evaluator
\item \emph{Dual resolution strategies}: a design that maximizes the affordances of two distinct execution environments---runtime reflection on the JVM and static lookup in GraalVM native image---rather than reducing both to a common denominator
\item \emph{Fail-soft error discipline}: appropriate for interactive creative coding, where partial results are more useful than hard failures
\end{itemize}

These contributions situate Bioscoop within a lineage of transpilers that target domain notation languages---Sass targeting CSS, Hiccup targeting HTML---where the value is not in replacing the target language but in providing the abstraction mechanisms it deliberately omits. Its most practically significant individual property is the \texttt{for} construct, which closes the expressiveness gap between FFmpeg's static filtergraph syntax and the data-driven, parameterized pipeline descriptions that users of complex filtergraphs actually need: graph topology and filter arguments derived from computed values, not written out by hand.
The system is available as open-source software under the MIT license.

\acks
The author thanks the Clojure community for their contributions to the ecosystem that made this work possible, and the FFmpeg project for building a framework expressive enough to warrant a compiler.

\printbibliography

\end{document}
